\chapter{Random Signals and Noise}
\label{ch:noise}

\begin{chapterintro}
If there were no noise, a single sample of a voltage could carry an infinite number of
bits, and communication engineering would be trivial. Noise is what makes the subject
hard and interesting. This chapter gives the probability and random-process tools the rest
of the book needs, explains where noise comes from physically, shows how engineers
characterise receivers with noise figure and noise temperature, and ends with the link
budget---the single calculation that decides whether a communication system will work at
all, from a Wi-Fi access point to a spacecraft at the edge of the solar system.
\end{chapterintro}

\section{The floor under everything}

In 1928 John B.~Johnson, measuring amplifiers at Bell Labs, found that every resistor
produced a small fluctuating voltage whose power was proportional to its temperature and
to the measurement bandwidth, and independent of what the resistor was made of. Harry
Nyquist explained the result the same year from thermodynamics. Thermal noise is not a
defect of any particular device; it is the random motion of charge carriers in any
conductor above absolute zero, and it sets a floor below which no receiver can hear. Every
other impairment---distortion, interference, fading---can in principle be designed away.
Noise can only be outwitted: by using more power, more bandwidth, more antennas, longer
observation, or cleverer codes. To reason about it we need probability.

\section{Random variables in communication}

\subsection{The Gaussian distribution and the central limit theorem}

A random variable $X$ is described by its probability density function (pdf)
$p_X(x)$, its mean $\mu=\E[X]$ and variance $\sigma^2=\E[(X-\mu)^2]$. The most important
distribution in communications is the Gaussian (normal),
\begin{equation}
p_X(x)=\frac{1}{\sqrt{2\pi\sigma^2}}\,e^{-(x-\mu)^2/2\sigma^2}.
\end{equation}
Its importance comes from the \term{central limit theorem}: the sum of many independent
random contributions of comparable size tends to a Gaussian, whatever the individual
distributions. Thermal noise is the sum of the effects of an astronomical number of
electrons; the interference from many distant transmitters, the error of a quantizer with
many levels and the sum of many multipath components all tend to Gaussian for the same
reason. Figure~\ref{fig:ch03_clt} shows how quickly the convergence happens: the sum of
just four uniform random variables is already hard to tell from a Gaussian by eye, though
the tails---which matter for error probability---converge more slowly.

\mdcfig{ch03_clt}{The central limit theorem in action: normalised sums of $n$ independent
uniform random variables (histograms) approach the Gaussian density (red curves).}

\subsection{The Q function}

Nearly every error probability in this book is a Gaussian tail probability. We define the
\term{Q function} as the probability that a zero-mean, unit-variance Gaussian exceeds $x$,
\begin{equation}
\Q(x)=\frac{1}{\sqrt{2\pi}}\int_x^\infty e^{-u^2/2}\,du=\tfrac12\operatorname{erfc}\!\left(x/\sqrt2\right).
\label{eq:ch03:q}
\end{equation}
If a receiver decides between $+A$ and $-A$ corrupted by Gaussian noise of standard
deviation $\sigma$, the error probability is $\Q(A/\sigma)$. The function falls extremely fast
(Figure~\ref{fig:ch03_q}): $\Q(3.09)=10^{-3}$, $\Q(4.75)=10^{-6}$, $\Q(6.0)=10^{-9}$.
Increasing the argument by about 50\%---a 3.5\,dB increase in signal-to-noise
ratio---reduces the error rate by a factor of a thousand. This steepness is why a working
digital link has a ``cliff'': a few decibels separate perfect reception from none. The
bounds $\Q(x)\le\frac12e^{-x^2/2}$ and $\Q(x)<e^{-x^2/2}/(x\sqrt{2\pi})$ are useful for quick
estimates and for proofs.

\mdcfig[0.8]{ch03_q}{The Q function (Gaussian tail probability) on a logarithmic scale, with
two classical bounds. Values at three reference error rates are marked.}

\subsection{Complex Gaussian noise and its envelope}

Because receivers work with complex baseband signals (Section~\ref{sec:complexenvelope}),
noise is complex: $n=n_I+jn_Q$ with independent zero-mean Gaussian parts of equal
variance $\sigma^2=N_0/2$. Such noise is \emph{circularly symmetric}: its distribution does not
change under rotation, so its phase is uniform on $[-\pi,\pi)$ and its total variance is
$\E|n|^2=N_0$. The magnitude $R=|n|$ has the \term{Rayleigh distribution}
\begin{equation}
p_R(r)=\frac{r}{\sigma^2}\,e^{-r^2/2\sigma^2},\qquad r\ge0,
\end{equation}
and the power $|n|^2$ is exponentially distributed. If a constant component $s$ is added---a
sinusoid plus noise, or a line-of-sight path plus scattered paths---the envelope follows the
\term{Rice distribution},
\begin{equation}
p_R(r)=\frac{r}{\sigma^2}\exp\!\left(-\frac{r^2+s^2}{2\sigma^2}\right)I_0\!\left(\frac{rs}{\sigma^2}\right),
\end{equation}
characterised by the \emph{K-factor} $K=s^2/2\sigma^2$, the ratio of the steady power to the
random power (Figure~\ref{fig:ch03_envelopes}). Stephen Rice derived these distributions in
1944 to describe radio noise; they returned as the models of multipath fading in
Chapter~\ref{ch:channels}.

\mdcfig{ch03_envelopes}{Left: envelope distributions of a constant plus circular Gaussian
noise, for three Rice K-factors (histograms from simulation, curves from theory).
Right: the phase of circularly symmetric Gaussian noise is uniformly distributed.}

\section{Random processes}

\subsection{Stationarity, correlation and spectrum}

A \term{random process} $x(t)$ is a random variable at every instant. We describe it by its
mean and its autocorrelation function $R_x(t_1,t_2)=\E[x(t_1)x^*(t_2)]$. If the mean is constant
and the autocorrelation depends only on the difference $\tau=t_1-t_2$, the process is
\emph{wide-sense stationary} (WSS), and its power spectral density is the Fourier transform
of the autocorrelation---the \term{Wiener--Khinchin theorem}:
\begin{equation}
S_x(f)=\int R_x(\tau)e^{-j2\pi f\tau}d\tau,\qquad \E|x(t)|^2=R_x(0)=\int S_x(f)\,df.
\end{equation}
The autocorrelation measures how quickly a process ``forgets'' its past; the PSD says how its
power is spread over frequency. A process that changes quickly has a narrow
autocorrelation and a wide spectrum, and vice versa---the same time--bandwidth reciprocity as
for deterministic signals.

Many communication signals are not quite stationary: a data stream sent at symbol rate
$R_s$ has statistics that repeat with period $1/R_s$. Such \emph{cyclostationary} processes
are handled by averaging the autocorrelation over one period, which yields the PSD
formulas of Chapter~\ref{ch:baseband}. Cyclostationarity is also a resource: it lets a
receiver detect and classify signals buried in noise, because noise has no such periodic
structure.

\subsection{Filtering random processes}

When a WSS process passes through an LTI filter $H(f)$, the output is WSS with
\begin{equation}
S_y(f)=|H(f)|^2S_x(f).
\label{eq:ch03:filtpsd}
\end{equation}
This single equation carries most of the noise analysis in this book. It says, for example,
that noise power at a receiver output is the input noise PSD integrated over the
receiver's \emph{power} response, which is why the noise-equivalent bandwidth of
Section~\ref{sec:bandwidth} is the right bandwidth to use.

\subsection{White noise}

Thermal noise has a flat PSD from DC to far beyond any radio frequency (up to about
6\,THz at room temperature, where quantum effects set in). We idealise it as \term{white
noise} with two-sided PSD $N_0/2$ W/Hz, whose autocorrelation is an impulse:
$R_n(\tau)=\frac{N_0}{2}\delta(\tau)$. White noise has infinite power and does not exist, but
once filtered it becomes perfectly well behaved: through a filter of noise bandwidth $B$,
the output power is $N_0B$ (counting both positive and negative frequencies of a real
filter). Figure~\ref{fig:ch03_filtered_noise} shows white noise and low-pass filtered noise:
filtering smooths the sample paths, narrows the spectrum and widens the autocorrelation.

\mdcfig{ch03_filtered_noise}{White Gaussian noise before and after a 50\,Hz low-pass filter:
sample paths, power spectral densities and autocorrelation functions.}

\subsection{Bandpass noise}

At the output of a radio's IF filter, noise occupies a band around the carrier. Its complex
envelope, like any bandpass signal's, can be written $\tilde{n}(t)=n_I(t)+jn_Q(t)$. If the
filter is symmetric about $f_c$, the in-phase and quadrature noise components are
independent, each with PSD $N_0$ over the baseband bandwidth $B/2$, and the complex
baseband noise has PSD $N_0$ across $|f|<B/2$. This is the noise model used in every
simulation in the labs: independent Gaussian I and Q samples, each of variance $N_0/2$.

\begin{pitfall}[Factors of two]
Noise PSDs are quoted one-sided ($N_0$ W/Hz over $f>0$) or two-sided ($N_0/2$ over all
$f$), and complex baseband noise has variance $N_0$ split equally between I and Q. Most
errors of exactly 3\,dB in simulations come from mixing these conventions. The labs use
one convention throughout: $\E|n|^2=N_0$ per complex sample and unit-energy symbols, so that
$E_s/N_0$ is the ratio of symbol energy to per-sample noise variance at the matched-filter
output.
\end{pitfall}

\section{Physical noise sources}

\subsection{Thermal noise}

A resistor $R$ at absolute temperature $T$ can deliver to a matched load an available noise
power
\begin{equation}
P_n=kTB,\qquad k=1.380649\times10^{-23}\ \mathrm{J/K},
\label{eq:ch03:ktb}
\end{equation}
independent of $R$. At the reference temperature $T_0=290$\,K, $kT_0=4.00\times10^{-21}$\,W/Hz,
or $-174$\,dBm/Hz. This number deserves memorising, together with its common scalings:
$-144$\,dBm in 1\,kHz, $-114$\,dBm in 1\,MHz, $-101$\,dBm in 20\,MHz. A GPS signal arrives at
about $-130$\,dBm spread over 2\,MHz, some 16\,dB \emph{below} the thermal noise in that band---yet
receivers track it routinely, by despreading (Chapter~\ref{ch:spreadspectrum}).

\subsection{Other noise mechanisms}

\emph{Shot noise} arises because current consists of discrete charges crossing a barrier at
random times; its PSD is $2qI$ (one-sided, in A$^2$/Hz), and it dominates in photodiodes and
therefore in optical receivers (Chapter~\ref{ch:wireline}). \emph{Flicker} or $1/f$ noise
rises at low frequencies in all active devices; it is why zero-IF receivers, which move the
signal to DC, must fight excess low-frequency noise, and why oscillators have close-in phase
noise that rises steeply near the carrier. At very high frequencies and low temperatures
the classical $kT$ gives way to the quantum limit $hf$; at 300\,GHz and 290\,K the two are
still within a factor of 20, so quantum noise is relevant only in optical and deep-cryogenic
systems.

\section{Noise figure and noise temperature}

\subsection{Definitions}

Every real amplifier adds noise. We characterise a two-port device by its \term{noise
factor} $F$, the ratio of the input SNR to the output SNR when the input noise is thermal at
$T_0=290$\,K; its \term{noise figure} is $\mathrm{NF}=10\log_{10}F$ dB. Equivalently, the device
can be modelled as noiseless with an extra noise source at its input of \term{equivalent
noise temperature}
\begin{equation}
T_e=(F-1)\,T_0 .
\end{equation}
A 1\,dB noise figure corresponds to $T_e=75$\,K; 3\,dB to 290\,K; 0.5\,dB to 35\,K. Noise
temperature is the natural unit when the input noise is not at 290\,K---for example, an
antenna looking at the cold sky---and it is the unit satellite and radio-astronomy engineers
prefer. A passive lossy component (cable, filter, attenuator) with loss $L$ at physical
temperature $T_0$ has noise factor $F=L$: a 2\,dB cable in front of a receiver adds exactly
2\,dB to the noise figure.

\subsection{Cascades: the Friis formula}

For a chain of stages with gains $G_i$ and noise factors $F_i$, Harald Friis showed in 1944
that
\begin{equation}
F=F_1+\frac{F_2-1}{G_1}+\frac{F_3-1}{G_1G_2}+\cdots,\qquad
T_e=T_{e1}+\frac{T_{e2}}{G_1}+\frac{T_{e3}}{G_1G_2}+\cdots
\label{eq:ch03:friis}
\end{equation}
The noise of each stage is divided by the total gain before it. The first stage dominates,
provided it has enough gain; this is the reason every sensitive receiver starts with a
low-noise amplifier (LNA), and why the LNA should be as close to the antenna as possible.
Figure~\ref{fig:ch03_friis} makes the point numerically: moving a 1\,dB LNA from after a 2\,dB
cable and a 1.5\,dB filter to directly at the antenna lowers the system noise figure from
about 5.1\,dB to 2.2\,dB, the equivalent of doubling the transmitter power.

\mdcfig{ch03_friis}{Contribution of each stage to the cascade noise factor for two orderings of
the same components. Losses in front of the LNA add directly; everything after a
high-gain LNA barely matters.}

\begin{inpractice}[Mast-head amplifiers and satellite LNBs]
Cellular base stations often place a tower-mounted amplifier at the top of the mast, before
tens of metres of lossy coaxial cable. A satellite TV dish goes further: its low-noise block
downconverter (LNB) at the focus amplifies the 10.7--12.75\,GHz signal with a noise
temperature under 100\,K and immediately converts it to 950--2150\,MHz, where cable losses
are tolerable. On the Ettus B200, the AD9364's internal LNA gives a noise figure under
8\,dB---fine for bench work, but an external LNA with 1\,dB NF will improve weak-signal
reception by several decibels if placed at the antenna.
\end{inpractice}

\subsection{Antenna and system noise temperature}

An antenna collects noise from whatever it sees. Pointed at the cold sky at microwave
frequencies, it may see only 10--50\,K; pointed at the ground, about 290\,K; at HF, galactic
and man-made noise can reach millions of kelvin (Figure~\ref{fig:ch03_skynoise}). The
\emph{system noise temperature} referred to the antenna terminals is
$T_{\text{sys}}=T_{\text{ant}}+T_e$, and the receiver's noise power is $kT_{\text{sys}}B$. The
figure explains two facts of radio engineering. At HF (3--30\,MHz), external noise is so
strong that receiver noise figure hardly matters: a 10\,dB NF receiver is as good as a
perfect one. In the ``microwave window'' from about 1 to 10\,GHz, the sky is quiet and
receiver noise dominates, which is why satellite, deep-space and radio-astronomy systems
live there and use cryogenically cooled amplifiers.

\mdcfig[0.88]{ch03_skynoise}{Illustrative external noise temperatures versus frequency.
Below about 30\,MHz, galactic and man-made noise dwarf receiver noise; between roughly 1 and
10\,GHz the sky is cold and receiver noise dominates; above 20\,GHz atmospheric absorption
adds noise. (Trends after ITU-R P.372; values are order-of-magnitude.)}

\section{SNR, $E_b/N_0$ and $C/N_0$}

Three ratios describe how strong a signal is relative to noise, and translating between
them is daily work.
\begin{itemize}
\item The \textbf{signal-to-noise ratio} $\mathrm{SNR}=P/(N_0B)$ depends on the bandwidth $B$
in which it is measured.
\item The \textbf{carrier-to-noise-density ratio} $C/N_0=P/N_0$, in dB-Hz, is independent of
bandwidth. It is the natural figure of merit for a link budget, because it can be computed
before the modulation is chosen. GPS receivers report it for each satellite; 45\,dB-Hz is a
strong GPS signal.
\item The \textbf{energy per bit to noise density} $E_b/N_0=P/(R_bN_0)$ normalises by the bit
rate $R_b$. It is the fair basis for comparing modulation and coding schemes, because it
charges each scheme for the energy it spends per bit of information.
\end{itemize}
They are related by
\begin{equation}
\frac{E_b}{N_0}=\frac{C}{N_0}\cdot\frac{1}{R_b}=\mathrm{SNR}\cdot\frac{B}{R_b},
\qquad
\frac{E_s}{N_0}=\frac{E_b}{N_0}\cdot\log_2M\cdot R_c ,
\label{eq:ch03:ebno}
\end{equation}
where $E_s$ is the energy per symbol of an $M$-ary scheme with code rate $R_c$.

\begin{worked}[From $C/N_0$ to bit rate]
A receiver measures $C/N_0=60$\,dB-Hz. A QPSK link with a rate-1/2 code needs
$E_b/N_0=2$\,dB at the target error rate. The maximum bit rate is
$R_b=(C/N_0)/(E_b/N_0)=10^{(60-2)/10}\approx630$\,kb/s. With 3\,dB of implementation margin,
plan for 315\,kb/s.
\end{worked}

\section{The link budget}

\subsection{The Friis transmission equation}

In free space, the power received by an antenna of gain $G_r$ at distance $d$ from a
transmitter of power $P_t$ and antenna gain $G_t$ is
\begin{equation}
P_r=P_tG_tG_r\left(\frac{\lambda}{4\pi d}\right)^2,
\qquad
L_{\text{fs}}=\left(\frac{4\pi d}{\lambda}\right)^2 .
\label{eq:ch03:friistx}
\end{equation}
The free-space path loss $L_{\text{fs}}$ appears to grow with frequency, but that is an artefact of
using antennas of fixed gain: an antenna of fixed \emph{area} $A$ has gain $4\pi A/\lambda^2$, so
between two dishes of fixed size the received power actually \emph{rises} with frequency. This
is why point-to-point microwave links and satellites use high frequencies and large dishes.
In decibels, the link budget is a sum:
\begin{equation}
P_r\,[\mathrm{dBm}]=P_t+G_t+G_r-L_{\text{fs}}-L_{\text{other}},\qquad
L_{\text{fs}}\,[\mathrm{dB}]=32.45+20\log_{10}f_{\mathrm{MHz}}+20\log_{10}d_{\mathrm{km}}.
\end{equation}

\subsection{Margins}

A real link budget adds losses for cables, polarisation mismatch, rain, walls, body
blockage and, above all, \emph{fading}. Multipath fading can drop the received power 20--30\,dB
below its average for a small fraction of the time (Chapter~\ref{ch:channels}), so the
budget must include a \emph{fade margin} chosen for the required availability. The final
comparison is between the received $E_b/N_0$ or SNR and the value the modem needs, plus
margin. Figure~\ref{fig:ch03_linkbudget} draws an indoor Wi-Fi budget as a waterfall: 20\,dBm
from the access point, antenna gains of 3\,dBi, about 102\,dB of path loss at 50\,m with a
path-loss exponent of 3 at 5.5\,GHz, and a wall, leaves about $-84$\,dBm---14\,dB above the noise
floor of a 20\,MHz receiver with 7\,dB noise figure. That supports 16-QAM with moderate
coding (Chapter~\ref{ch:wifi}).

\mdcfig{ch03_linkbudget}{A link budget drawn as a waterfall. Gains lift the level, losses lower it;
the gap between the received power and the noise floor is the SNR available to the modem.}

\begin{history}[Talking to Voyager]
Voyager~1, launched in 1977, is now more than 160 times as far from the Sun as the Earth
is. It transmits roughly 20\,W at 8.4\,GHz from a 3.7\,m dish, and its signal is received by
the 70\,m antennas of NASA's Deep Space Network, whose cryogenic receivers have system noise
temperatures of a few tens of kelvin. At more than $2.4\times10^{10}$\,km, the free-space
path loss exceeds 300\,dB, and the received power is below $10^{-19}$\,W. The link still
carries 160\,b/s of science data, thanks to high antenna gains, very low noise temperatures,
a narrow bandwidth and powerful concatenated error-correcting codes---a convolutional code
and a Reed--Solomon code, the same pair of codes that protected digital television for a
generation (Chapter~\ref{ch:classiccodes}). Every term of Equation~\eqref{eq:ch03:friistx}
has been pushed to its limit.
\end{history}

\begin{labbox}[Noise and link budgets]
Lab~1 contains an interactive sensitivity calculator: vary bandwidth, noise figure and
required SNR and watch the free-space range change. The figure script
\texttt{figscripts/ch03\_figs.py} contains the cascade-noise-figure calculator used for
Figure~\ref{fig:ch03_friis}; rearrange the stages of your own receiver and find the best
placement for the filter. Chapter~\ref{ch:modulation}'s lab will convert these SNRs into
bit error rates.
\end{labbox}

\problems
\begin{enumerate}[label=\thechapter.\arabic*]
\item Show that if $X$ and $Y$ are independent $\mathcal{N}(0,\sigma^2)$, then
$R=\sqrt{X^2+Y^2}$ is Rayleigh distributed and $R^2$ is exponential. What fraction of the time
is a Rayleigh-fading signal more than 10\,dB below its mean power? 20\,dB?
\item Compute the thermal noise power in dBm for bandwidths of 12.5\,kHz (narrowband FM),
200\,kHz (GSM), 20\,MHz (Wi-Fi) and 400\,MHz (5G mmWave), all at 290\,K.
\item A receiver has an antenna with $T_{\text{ant}}=50$\,K, a feed line with 0.5\,dB loss at
290\,K, an LNA with NF 0.8\,dB and 30\,dB gain, and a second stage with NF 10\,dB. Compute the
system noise temperature. How much does moving the LNA before the feed line improve it?
\item Derive the Friis formula~\eqref{eq:ch03:friis} for two stages from the definition of
noise temperature.
\item A geostationary satellite 36\,000\,km away transmits 50\,dBW EIRP at 12\,GHz. A 60\,cm
dish with 65\% efficiency and a 60\,K system noise temperature receives it. Compute $C/N_0$.
What bit rate can a DVB-S2 QPSK rate-3/4 carrier (needing about 4\,dB $E_s/N_0$) support?
\item A Bluetooth Low Energy link transmits 0\,dBm at 2.4\,GHz with 0\,dBi antennas. The
receiver sensitivity is $-95$\,dBm. What is the free-space range? Repeat with a path-loss
exponent of 3 beyond 1\,m.
\item \emph{(Simulation)} Generate complex Gaussian noise in Python and verify numerically that
(a) its phase is uniform, (b) its magnitude is Rayleigh, and (c) after a filter $H(f)$ its
PSD is $|H(f)|^2N_0$. Then measure how many samples are needed to estimate a $10^{-6}$ tail
probability to within 10\%.
\end{enumerate}

\furtherreading
\begin{itemize}
\item A.~Papoulis and S.~U.~Pillai, \emph{Probability, Random Variables and Stochastic
Processes}, 4th ed., McGraw-Hill, 2002.
\item S.~O.~Rice, ``Mathematical analysis of random noise,'' \emph{Bell System Technical
Journal}, 1944--45. The classic.
\item H.~T.~Friis, ``Noise figures of radio receivers,'' \emph{Proc. IRE}, 32, 1944.
\item Recommendation ITU-R P.372, \emph{Radio noise}. Measured and modelled external noise
levels from 0.1\,Hz to 100\,GHz.
\item J.~H.~Yuen (ed.), \emph{Deep Space Telecommunications Systems Engineering}, JPL,
1982 (free online). The engineering of links like Voyager's.
\end{itemize}
